On considère $7$ candidats d'un fokontany :
\[
\mathcal{E} = \{A, B, C, D, E, F, G\}.
\]
On souhaite former un comité d'exactement $3$ personnes.

À chaque comité, on associe un mot binaire de longueur $7$, dans lequel $1$
signifie que le candidat est choisi et $0$ qu'il ne l'est pas.

Par exemple,
\[
\{A, D, G\} \longleftrightarrow 1001001.
\]
\begin{enumerate}
  \item Quel comité correspond au mot
  \[
  0110100 ?
  \]
  \item Quel mot correspond au comité
  \[
  \{B, E, F\} ?
  \]
  \item Combien de chiffres $1$ possède le mot associé à un comité de $3$
  personnes ?
  \item Montrer que les comités de $3$ personnes sont en bijection avec les
  mots binaires de longueur $7$ contenant exactement $3$ chiffres égaux à $1$.
  \item En déduire que le nombre de comités est
  \[
  \binom{7}{3}.
  \]
  \item Généraliser : avec $n$ candidats, combien existe-t-il de comités de
  $k$ personnes ?
\end{enumerate}
